\section{Methods}

\subsection{Problem formulation}

We study drug-target affinity prediction under an anchor-transfer formulation. Instead of scoring a protein-drug pair in isolation, the model receives an anchor protein $a$, a query protein $q$, and a drug $d$, and predicts both a continuous affinity value and a binary binding score:
\[
f(a,q,d)\rightarrow (\hat{y},\hat{b}),
\]
where $\hat{y}$ is the predicted pKi and $\hat{b}\in[0,1]$ is the predicted binding probability. The central idea is that the anchor provides binding-specific context for the query drug. If a drug is already known to bind one protein strongly, the model can ask whether the query protein is compatible with that same binding context. This formulation draws on the broader principle of transfer learning \cite{pan2010transfer}, in which knowledge from a source domain (the known anchor interaction) is reused to improve prediction in a target domain (the unseen query interaction).

For each drug in the training set, we define the anchor as the protein with the highest observed pKi for that drug. Thus, given a query pair $(q,d)$, the corresponding anchor $a(d)$ is the strongest known binder of $d$ in the source data. When the query protein itself is the top-ranked binder, we use the next-highest protein so that the anchor and query are not identical. This produces training triplets of the form $(a,q,d)$ with supervision on the affinity of the query pair $(q,d)$.

We train the model jointly for regression and classification, following a multitask learning approach \cite{caruana1997multitask}. Continuous affinity supervision uses the observed pKi value $y$, while the binary target is derived by thresholding pKi into confident binders and non-binders. We assign $b=1$ when $y\geq 7$, $b=0$ when $y\leq 5$, and treat the intermediate range $5<y<7$ as ambiguous. These samples remain useful for regression but are excluded from the binary loss. This multitask setup lets the model learn both fine-grained affinity ranking and coarse binding discrimination.

\subsection{Anchor transfer architecture}

Figure~\ref{fig:architecture} summarizes the proposed model. The architecture contains a shared protein encoder, a shared drug encoder, three pairwise interaction multilayer perceptrons, and two prediction heads. The shared encoders produce aligned 256-dimensional embeddings for the anchor protein, query protein and drug; the interaction blocks then model all pairwise relations inside the triplet.

\subsection{Protein encoder}

Each protein is represented by a frozen embedding from ESM-2 \cite{lin2023esm2}. We use either the 35M model, which produces a 480-dimensional sequence embedding, or the 650M model, which produces a 1,280-dimensional embedding. The pretrained ESM-2 parameters are kept fixed throughout training, consistent with the common practice of using frozen protein language model features for downstream prediction tasks \cite{rives2021biological,rao2019evaluating}. To map both anchor and query proteins into the same task-specific space, we apply the same projection module to both branches:
\[
h_p = \mathrm{LayerNorm}\!\left(\mathrm{ReLU}(W_p x_p + b_p)\right),
\]
where $x_p\in\mathbb{R}^{480}$ or $\mathbb{R}^{1280}$ is the frozen ESM-2 representation, $h_p\in\mathbb{R}^{256}$ is the projected embedding, and LayerNorm denotes layer normalization \cite{ba2016layernorm}. Weight sharing is important here: anchor and query proteins must be directly comparable, so they are encoded with the same linear projection, nonlinearity and normalization pipeline.

\subsection{Drug encoder}

Drugs are encoded from SMILES strings \cite{weininger1988smiles} using a compact convolutional encoder inspired by DeepDTA \cite{ozturk2018deepdta}. After token embedding, the SMILES sequence is processed by three parallel one-dimensional convolution branches with kernel sizes 4, 6 and 8. Each branch uses 32 filters followed by global max pooling. The pooled branch outputs are concatenated into a 96-dimensional drug feature,
\[
z_d = \big[\mathrm{pool}(\mathrm{Conv}_4(d)) \,\|\, \mathrm{pool}(\mathrm{Conv}_6(d)) \,\|\, \mathrm{pool}(\mathrm{Conv}_8(d))\big]\in\mathbb{R}^{96},
\]
and a final linear projection maps this vector to the shared 256-dimensional latent space:
\[
h_d = W_d z_d + b_d.
\]
Using parallel kernels allows the encoder to capture SMILES motifs of different lengths while keeping the drug branch lightweight.

\subsection{Triple pairwise interaction module}

Once the three entities are encoded, the model explicitly constructs interaction features for every pair in the triplet. Let $h_a$, $h_q$ and $h_d$ denote the 256-dimensional embeddings of the anchor protein, query protein and drug. We then compute
\[
h_{ad} = \phi_{ad}([h_a \| h_d]), \qquad
h_{qd} = \phi_{qd}([h_q \| h_d]), \qquad
h_{aq} = \phi_{aq}([h_a \| h_q]),
\]
where each $\phi$ is an independent pairwise multilayer perceptron that takes the 512-dimensional concatenated input and projects it to 256 dimensions followed by a ReLU nonlinearity. These three learned features capture distinct aspects of the problem: the anchor-drug branch represents evidence that the drug already binds a known reference protein, the query-drug branch represents direct compatibility between the query and the compound, and the anchor-query branch represents similarity or complementarity between the reference protein and the query.

\subsection{Fusion and prediction heads}

The final representation concatenates the three base embeddings and the three pairwise features:
\[
z = [h_a \| h_q \| h_d \| h_{ad} \| h_{qd} \| h_{aq}] \in \mathbb{R}^{1536}.
\]
This fused vector is passed to two task-specific heads. The classifier predicts binding probability through a multilayer perceptron followed by a sigmoid output, while the regressor predicts continuous pKi with the same hidden structure but a linear output. The classifier is optimized for binary binding discrimination, whereas the regressor preserves the quantitative affinity signal. In practice, the dual-head design is useful because AUROC and pKi regression quality do not always move together under dataset shift.

\subsection{Training objective}

The overall loss combines binary cross-entropy and mean squared error:
\[
\mathcal{L} = \mathcal{L}_{\mathrm{BCE}} + \alpha \mathcal{L}_{\mathrm{MSE}}.
\]
In all experiments, we set $\alpha=1.0$.
The regression loss is computed over all examples,
\[
\mathcal{L}_{\mathrm{MSE}} = \frac{1}{N}\sum_{i=1}^{N}(\hat{y}_i-y_i)^2,
\]
whereas the classification loss is masked so that ambiguous affinities do not contribute. Defining
\[
m_i=
\begin{cases}
1, & y_i\leq 5 \text{ or } y_i\geq 7, \\
0, & 5<y_i<7,
\end{cases}
\]
we compute
\[
\mathcal{L}_{\mathrm{BCE}} =
\frac{1}{\sum_i m_i}
\sum_{i=1}^{N}
m_i\,\mathrm{BCE}(\hat{b}_i,b_i).
\]
This masking scheme preserves the full pKi range for regression while avoiding noisy supervision near the decision boundary. It also matches the evaluation setting, where binary discrimination is most meaningful for clearly binding and clearly non-binding interactions.

\subsection{Drug-anchor ablation}

To test whether the gains come from protein-side anchoring specifically or from conditioning on any known same-partner example, we evaluate a drug-anchor ablation. In this variant, the inputs are an anchor drug $a_d$, a query drug $q_d$, and a protein $p$, and the model predicts pKi for the query pair:
\[
f(a_d,q_d,p)\rightarrow \hat{y}.
\]
The protein branch remains the frozen ESM-2 projection, while the anchor-drug and query-drug branches share the same SMILES convolutional encoder. Pairwise interaction blocks are built for $(a_d,p)$, $(q_d,p)$ and $(a_d,q_d)$, mirroring the structure of the main model. This ablation tests whether transferable context is still useful when the anchor is moved to the compound side rather than the protein side.

\subsection{ESM-DTA baseline}

As a pairwise baseline, we construct \emph{ESM-DTA} by replacing DeepDTA's learned protein character embedding pathway with frozen ESM-2 features \cite{lin2023esm2}. The drug encoder remains the same DeepDTA-style SMILES convolution used in the main model, but the input is only the pair $(q,d)$ and there is no anchor. Protein and drug embeddings are concatenated and passed to a standard multilayer perceptron to predict affinity. This baseline is intentionally simple: it asks whether stronger pretrained protein representations alone are sufficient, or whether the improvement instead comes from explicitly conditioning prediction on a known strong binder for the same drug.

\subsection{Conformation weighting analysis}

As an auxiliary analysis for IDP queries, we replaced the query protein's ESM-2 embedding with a TM-weighted average of ESM-2 embeddings from structurally matched ordered proteins. The matches were obtained by running Foldseek on IDRome conformations and aggregating the retrieved ordered-protein embeddings by TM-score, while leaving the rest of the model and evaluation protocol unchanged.
